\paragraph{\trFormalDiffeomorphism。}\label{sub:formal-diffeomorphisms}
这里给出（某一类特殊的）\trAFunctor{}的一个稍有不同的视角。设我们有一个由\trObject{}组成的
集合 $\Ob\scr{A}$，并且对每一对\trObject{}都配有\trGraded\trVectorSpace{}%
$\hom_{\scr{A}}(X_0,X_1)$，只是暂时不设任何\trComposition{}映射。所谓从 $\scr{A}$ 到自身的
\textbf{\trFormalDiffeomorphism}（formal diffeomorphism）$\Phi$，就是任意一列形如
\eqref{eq:afunctor-components} 的映射，唯一的条件是：每个 $\Phi^1$ 都应当是
$\hom_{\scr{A}}(X_0,X_1)$ 的\trLinear\trAutomorphism。公式\eqref{eq:afunctor-composition}
使\trFormalDiffeomorphism{}成为一个群。现在再设我们有一组映射 $\mu^{\bullet}_{\scr{A}}$，
它们使 $\scr{A}$ 成为一个\trNonunital{}\trACategory。给定 $\Phi$，通过递归地求解
\eqref{eq:afunctor-equations} 可以证明：存在另一个唯一的\trAinfty{}结构
$\mu'^{\bullet}_{\scr{A}}$（底\trVectorSpace{}相同），使得
$\Phi\colon\scr{A}\to\tilde{\scr{A}}$ 是一个\trNonunital{}\trAFunctor。这就定义了%
\trFormalDiffeomorphism{}群在所有可能的 $\mu^{\bullet}_{\scr{A}}$ 之集合上的一个作用；我们记
$\tilde{\scr{A}}=\Phi_{*}\scr{A}$。

\begin{remark}\label{thm:formal-diffeo-name}
名称“\trFormalDiffeomorphism”来自非交换几何：在那里，人们可以把\trAinfty{}结构想象成%
\trGraded{}非交换仿射空间上\trDegree{}为 $-1$ 的\trVectorField，而这个作用不过是坐标变换
罢了。
\end{remark}
